\section{Objectives under Study}
\label{sec:method}

\paragraph{Hard-negative contrastive loss.}
A batch contains $B$ groups, hence $2B$ images and $2B$ captions. With $S$ the $2B{\times}2B$
score matrix, the loss $\mathcal{L}_{\mathrm{hn}}$ is the average of row-wise (image$\to$caption)
and column-wise (caption$\to$image) cross-entropies with diagonal targets. Every off-diagonal pair
that the oracle marks as true is removed from the denominators, so it is never used as a negative.
The group partner of each item is kept as a negative (the hard negative). This is the NegCLIP-style
use of minimal-change partners \citep{yuksekgonul2023bow,awal2024vismin}, applied here to a small
head on frozen features.

\begin{remark}[a standard bilinear identity]
\label{rem:identity}
Assume the score is an inner product $s(I,T)=\langle \phi(I),\psi(T)\rangle$, where $\phi$ depends
only on the image and $\psi$ only on the caption. Per-input normalization, as in cosine scores, is
allowed; normalization or attention that depends on the pair is not. Then, for a $2{\times}2$
group,
\begin{equation}
\begin{aligned}
S_{00}+S_{11}-S_{01}-S_{10}
&=\langle\phi(I_0),\psi(T_0)\rangle+\langle\phi(I_1),\psi(T_1)\rangle
-\langle\phi(I_0),\psi(T_1)\rangle-\langle\phi(I_1),\psi(T_0)\rangle\\
&=\langle \phi(I_0)-\phi(I_1),\,\psi(T_0)-\psi(T_1)\rangle ,
\end{aligned}
\label{eq:identity}
\end{equation}
by bilinearity alone. Both the head of Section~\ref{sec:setting} and zero-shot cosine CLIP satisfy
the assumption.
\end{remark}

Equation~\eqref{eq:identity} is elementary; we state it because it fixes what several signals
share. GroupMatch is exactly the event that the right-hand side is positive. The two within-group
row margins of the hard-negative loss, $S_{00}-S_{01}$ and $S_{11}-S_{10}$, sum to the same
quantity. Aligning the image edit vector $\phi(I_1)-\phi(I_0)$ with the caption edit vector
$\psi(T_1)-\psi(T_0)$ of the same group rewards the same inner product, up to normalization. None
of this holds automatically for scorers with pair-dependent computation, such as cross-attention
heads, and we do not apply it to them.

\paragraph{The edit-consistency objective under study.}
For each group $g$ in a batch let $\Delta^I_g=e(I^g_1)-e(I^g_0)$ and $\Delta^T_g=e(T^g_1)-e(T^g_0)$,
using the same member order in both modalities, where $e$ is the concatenated head output. The
objective $\mathcal{L}_{\mathrm{edit}}$ is a symmetric cross-entropy over the $B{\times}B$ matrix
$C_{gh}=\cos(\Delta^I_g,\Delta^T_h)/\tau$ with diagonal targets and $\tau=0.1$. Norms include a
small $\epsilon$, so zero edit vectors give uniform logits and a loss of $\log B$: unlike a squared
error between edit vectors, zero is not a minimizer. The diagonal entries reward the normalized
within-group quantity of Remark~\ref{rem:identity} and therefore overlap with the hard-negative
loss. The off-diagonal entries are the additional signal: an image edit must be closer in direction
to its own caption edit than to the caption edits of other groups in the batch. The objective
assumes that different groups in a batch carry different edits; edit-level false negatives are not
masked. The combined loss is $\mathcal{L}_{\mathrm{hn}}+\lambda\,\mathcal{L}_{\mathrm{edit}}$.

\paragraph{Scale.}
On the first training batch of the small panel, before training and with $\lambda=1$, the ratio
$\|\nabla\mathcal{L}_{\mathrm{edit}}\|/\|\nabla\mathcal{L}_{\mathrm{hn}}\|$ was \SPgradInitMin--\SPgradInitMax{}
across seeds; after hard-negative training it was \SPgradEndMin--\SPgradEndMax. We report these
ratios only to show that $\lambda$ must be small for the two terms to be on comparable scales. A
large gradient says nothing about usefulness.
